% Part: first-order-logic
% Chapter: axiomatic-deduction
% Section: axioms-rules-quantifiers

\documentclass[../../../include/open-logic-section]{subfiles}

\begin{document}

\olfileid{fol}{axd}{qua}

\olsection{بديهيات المكمّمات وقواعدها}

\begin{defn}[بديهيات المكمّمات]
\emph{البديهيات} التي تحكم المكمّمات هي جميع حالات الصيغ الآتية:
\begin{align}
\ollabel{ax:q1} & \lforall[x][!B] \lif !B(t), \\
\ollabel{ax:q2} & !B(t) \lif \lexists[x][!B],
\end{align}
لكل حد مغلق~$t$.
\end{defn}

% تصحيح مصدر (0118-I1): أُضيف مخزون محافظ غير متناهٍ للثوابت المميزة.
في الاشتقاقات الكمية نلحق بلغة المسألة~$\Lang L$ مجموعةً لا نهائية
$\mathcal E=\{e_0,e_1,\dots\}$ من الثوابت المساعدة الجديدة. ويجوز أن
تظهر ثوابت~$\mathcal E$ في السطور الوسطى، وتُقرأ مخططات البديهيات والقواعد
في اللغة الموسعة، على ألا يظهر شيء منها في دعم ال!!{derivation} أو خاتمته.
ونختار الثابت المميّز في كل تطبيق من~$\mathcal E$.

\begin{defn}[قواعد المكمّمات]
% تصحيح مصدر (0115-I1): أُضيف شرط عدم ورود الثابت المميز في A(x)، وهو ساقط من الأصل.
\begin{enumerate}
\item إذا ظهرت $!B \lif !A(a)$ بالفعل في ال!!{derivation} ولم يظهر
  $a$ في دعم ال!!{derivation} ولا في~$!B$ ولا في~$!A(x)$، فإن $!B \lif
  \lforall[x][!A(x)]$ خطوة استدلال صحيحة.
\item إذا ظهرت $!A(a) \lif !B$ بالفعل في ال!!{derivation} ولم يظهر
  $a$ في دعم ال!!{derivation} ولا في~$!B$ ولا في~$!A(x)$، فإن $\lexists[x][!A(x)] \lif
  !B$ خطوة استدلال صحيحة.
\end{enumerate}
\end{defn}

سنختصر كلًا من هاتين القاعدتين بالرمز «\QR».

ونعد من الثوابت الواقعة في شجرة برهان كل ثابت يظهر في صيغة فيها، وكل
ثابت مميز مكتوب على عقدة~\QR{}، ولو كان المتغير المسوّر لا يقع في الصيغة
فلم يظهر ذلك الثابت في مقدمة العقدة.

\begin{lem}[تجديد الثوابت المميزة]\ollabel{lem:qr-freshening}
لتكن~$\Pi$ !!a{derivation} ذات دعم منتهٍ وخاتمة~$!C$، ولتكن~$F$
مجموعة منتهية من الثوابت. توجد !!a{derivation}~$\Pi'$ خاتمتها~$!C$،
ودعمها محتوى في دعم~$\Pi$، وتكون الثوابت المميزة في تطبيقات~\QR{} فيها
متمايزة، ولا يقع شيء منها في~$F$ ولا في دعم~$\Pi$ ولا في~$!B$ أو
~$!A(x)$ الخاصين بتطبيقها.
\end{lem}

\begin{proof}
سمّ دعم~$\Pi$ الأصلي~$\Sigma_0$. احذف السطور غير المستعملة، وابسط
المتتالية الباقية إلى شجرة برهان منتهية دعمها~$\Sigma\subseteq\Sigma_0$.
نثبت بالاستقراء العبارة الأقوى الآتية. لكل شجرة منتهية~$T$ جذرها~$!D$
ودعمها~$\Sigma$، ولكل مجموعة منتهية~$H$ من الثوابت، توجد شجرة~$T'$
لها الجذر والدعم نفسيهما؛ وثوابتها المميزة متمايزة، ولا يقع شيء منها في~$H$ ولا في
$!D$ ولا في~$\Sigma$. وتظل شروط~\QR{} المحلية كلها متحققة. وتبقى
أوراق الفرض والبديهية على حالها.

في عقدة MP، لتكن شجرتا المقدمتين~$T_1,T_2$. جدّد~$T_1$ بعد أن تضم إلى
$H$ جميع الثوابت الواقعة في~$T_2$، ثم جدّد~$T_2$ بعد أن تضم إلى~$H$
جميع الثوابت الواقعة في الشجرة~$T'_1$ الناتجة. وهكذا لا يساوي ثابت مميز
في إحدى الشجرتين ثابتًا في الأخرى، ولا يقع في دعمها؛ فيبقى تطبيق MP نفسه
صحيحًا وتبقى شروط الجدة صحيحة بالنسبة إلى الدعم المجموع.

وعند عقدة~\QR{} ذات الثابت المميز~$a$ والجذر~$!D$، جدّد أولًا شجرة
المقدمة مع مجموعة التحريم~$H$ مضافًا إليها~$a$ وجميع ثوابت~$!D$.
ثم اختر~$b\in\mathcal E$ لا يقع في~$H$ ولا في شجرة المقدمة المجددة
كلها ولا في~$!D$، وبدّل الاسمين~$a$ و~$b$ بانتظام في شجرة المقدمة.
فالثابت~$a$ لا يقع، بحسب شرط التطبيق الأصلي، في الدعم ولا في~$!B$ ولا
في~$!A(x)$، والثابت~$b$ لم يكن واقعًا في الشجرة؛ لذلك لا يتغير الدعم،
ويتحول جذر المقدمة من الحالة ذات~$a$ إلى الحالة ذات~$b$، وتبقى خاتمة
عقدة~\QR{} هي~$!D$. وهذا التبديل تقابل في أسماء الثوابت؛ فينقل حالات
البديهيات إلى حالات بديهيات، ويحفظ MP والحدود المغلقة وعقد~\QR{}
الداخلية. كما أن الاستقراء حظر~$a$ على ثوابتها المميزة، واختيار~$b$
حظره على الشجرة كلها، فلا يندمج الثابت الجديد بأي ثابت مميز داخلي.
نعيد إذن تطبيق~\QR{} بالثابت~$b$. والحجة واحدة لصورتَي القاعدة.

وبهذا تتم العبارة الأقوى. نأخذ فيها~$H$ اتحاد~$F$ ومجموعة جميع الثوابت
الواقعة في~$\Sigma_0$، ثم نعيد الشجرة إلى متتالية منتهية. فيكون دعمها
$\Sigma\subseteq\Sigma_0$، ونتحصل على~$\Pi'$ المطلوبة.
\end{proof}

\end{document}
